\documentclass[12pt]{article}
\usepackage[utf8]{inputenc}
\usepackage[T1]{fontenc}
\usepackage[margin=1in]{geometry}
\usepackage{setspace}
\usepackage{xcolor}
\usepackage[colorlinks=true, urlcolor=blue!60!black, linkcolor=blue!60!black]{hyperref}
\setlength{\parindent}{2em}
\setlength{\parskip}{1.25em}
\renewcommand{\baselinestretch}{1.0}
\begin{document}
\begin{center}
{\LARGE \textit{Friends of the Earth, Inc. v. Laidlaw Environmental Services (TOC), Inc.}, 528 U.S. 167 (2000)\par}
\vspace{0.8em}
{\large Case Brief\par}
\end{center}
\vspace{1em}

\section{Memory Jogger}
Environmental groups sued a polluter under the Clean Water Act for civil penalties; the Court held 7--2 that members who curtailed use of a river had standing even without proof of environmental harm, that penalties deterring ongoing violations provide redress, and that voluntary cessation moots a case only if the defendant shows it is absolutely clear the violations cannot reasonably be expected to recur (\href{https://supreme.justia.com/cases/federal/us/528/167/}{Justia}).

\section{Detailed Case Facts}
In 1986, respondent Laidlaw Environmental Services (TOC), Inc. purchased a hazardous-waste incinerator facility in Roebuck, South Carolina, that included a wastewater treatment plant discharging into the North Tyger River. In 1987, the South Carolina Department of Health and Environmental Control (DHEC) issued Laidlaw a National Pollutant Discharge Elimination System (NPDES) permit under the Clean Water Act restricting discharges of pollutants, including mercury. Laidlaw repeatedly violated its permit's mercury limits---committing nearly 500 discharge violations between 1987 and 1995.

\section{Procedural History}
On April 10, 1992, petitioners Friends of the Earth (FOE) and Citizens Local Environmental Action Network (CLEAN)---later joined by the Sierra Club---sent Laidlaw the mandatory 60-day notice of intent to file a citizen suit under CWA \S~505(a), seeking declaratory and injunctive relief, civil penalties payable to the U.S. Treasury, and attorneys' fees. One day before the 60-day window expired, at Laidlaw's request, DHEC filed a state-court lawsuit against Laidlaw (Laidlaw itself drafted the complaint, paid the filing fee, and settled the next day for \$100,000). The federal District Court rejected Laidlaw's attempt to bar the citizen suit on ``diligent prosecution'' grounds, finding the state settlement hasty and inadequate to strip Laidlaw of its economic benefit from noncompliance.

Following trial, the District Court found that Laidlaw's mercury violations had caused ``no demonstrated proof of harm to the environment'' and denied injunctive relief because Laidlaw had achieved substantial compliance after the suit was filed, but assessed a \$405,800 civil penalty to deter future violations. FOE appealed the penalty amount as too low (without appealing the denial of injunctive relief), and Laidlaw cross-appealed on standing. Relying on \textit{Steel Co. v. Citizens for a Better Environment}, 523 U.S. 83 (1998), the Fourth Circuit vacated the penalty and ordered dismissal as moot (149 F.3d 303 (1998)), reasoning that because FOE sought only civil penalties payable to the Treasury on appeal, those penalties could not redress FOE's injury. While appellate proceedings were pending, Laidlaw permanently closed and dismantled the Roebuck incinerator, though it retained its NPDES permit.

\section{Judicial Votes}
The Supreme Court \textbf{reversed and remanded} the Fourth Circuit by a \textbf{7--2} vote. The case was argued on October 12, 1999, and decided on January 12, 2000 (\href{https://supreme.justia.com/cases/federal/us/528/167/}{Justia}; \href{https://www.law.cornell.edu/supremecourt/text/528/167}{Cornell LII}):

\begin{itemize}
    \item \textbf{Majority Opinion:} Justice Ginsburg delivered the opinion of the Court, joined by Chief Justice Rehnquist and Justices Stevens, O'Connor, Kennedy, Souter, and Breyer.
    \item \textbf{Concurring Opinions:} Justice Stevens filed a concurring opinion. Justice Kennedy filed a concurring opinion.
    \item \textbf{Dissenting Opinion:} Justice Scalia filed a dissenting opinion, joined by Justice Thomas.
\end{itemize}

\section{Holding}
\textbf{Questions Presented:}

\begin{enumerate}
  \item Did environmental organizations establish Article III standing---both injury in fact and redressability---to bring a Clean Water Act citizen suit seeking civil penalties payable to the U.S. Treasury for ongoing permit violations?
  \item Did Laidlaw's post-filing compliance with its permit or subsequent closure of its Roebuck facility render the citizen suit moot?
\end{enumerate}

\textbf{Holding:} Yes on standing; no on automatic mootness. Reversing the Fourth Circuit, the Supreme Court held that: (1) reasonable aesthetic, recreational, and economic concerns of members who curtailed their use of the North Tyger River satisfied Article III injury in fact, even without proof of physical harm to the river; (2) civil penalties payable to the U.S. Treasury redress ongoing violations at the time of suit by deterring future noncompliance; and (3) a defendant's voluntary cessation of unlawful conduct does not moot a case unless the defendant carries the ``formidable burden'' of making it ``absolutely clear'' that the wrongful behavior cannot reasonably be expected to recur---a distinct inquiry from standing that remained open for factual resolution on remand.

\section{Analysis of Opinions}
\subsection{Majority Opinion (Justice Ginsburg)}

Writing for a seven-Justice majority, Justice Ginsburg separated the Fourth Circuit's conflated analysis into two distinct Article III inquiries: standing (measured at the outset of litigation) and mootness (measured by intervening events).

\textbf{1. Injury in Fact (Injury to the Plaintiff, Not the Environment):} Under \textit{Hunt v. Washington State Apple Advertising Comm'n}, 432 U.S. 333 (1977), an association has standing if its members would have standing in their own right. Laidlaw relied on the District Court's penalty-phase finding of ``no demonstrated proof of harm to the environment'' to argue FOE suffered no injury. Justice Ginsburg squarely rejected that premise: the relevant Article III showing is ``not injury to the environment but injury to the plaintiff.'' Unlike the speculative ``some day'' intentions rejected in \textit{Lujan v. Defenders of Wildlife}, 504 U.S. 555 (1992), FOE's members submitted concrete affidavits showing they lived near the North Tyger River and had ceased fishing, camping, wading, boating, and picnicking because of Laidlaw's ongoing illegal mercury discharges. Because a reasonable fear of continuous unlawful pollution directly lessened the aesthetic and recreational value of the area, FOE established concrete injury in fact.

\textbf{2. Redressability (Deterrence via Public Civil Penalties):} Although a plaintiff must demonstrate standing separately for each form of relief, the Court held that civil penalties payable to the U.S. Treasury satisfy redressability when violations are \textit{ongoing or imminent at the time suit is filed}. Unlike \textit{Steel Co.}---where the defendant's violations were wholly past before the complaint was filed---civil penalties for ongoing violations serve a congressional deterrent purpose, encouraging violators to abate current discharges and discouraging recurrence. Thus, deterrence itself redresses a threatened plaintiff's injury.

\textbf{3. Mootness vs. Standing (Voluntary Cessation):} The Fourth Circuit erred by treating mootness as merely ``standing set in a time frame'' and using \textit{Steel Co.} to dismiss the civil-penalty appeal. Justice Ginsburg explained that standing and mootness differ critically in two respects: (a) \textit{Burden of Proof:} While the plaintiff bears the burden to establish standing at filing, a defendant claiming voluntary cessation moots a case bears the ``heavy'' and ``formidable'' burden of showing it is ``absolutely clear'' the challenged conduct cannot reasonably be expected to recur; and (b) \textit{Sunk Judicial Costs:} Exceptions like voluntary cessation and ``capable of repetition, yet evading review'' exist in mootness (but not standing) to prevent defendants from strategically evading judicial review after years of litigation. Because the lower courts had never resolved disputed facts about whether Laidlaw's compliance or plant closure (while retaining its permit) made recurrence impossible, the Court remanded.

\subsection{Concurring Opinion (Justice Stevens)}

Justice Stevens wrote separately to emphasize that post-filing voluntary compliance does not moot a claim for civil penalties already incurred for ongoing violations, because civil penalties serve punitive and deterrent functions distinct from prospective injunctive relief.

\subsection{Concurring Opinion (Justice Kennedy)}

Joining the Court's Article III analysis, Justice Kennedy added a terse separation-of-powers reservation: he noted that authorizing private citizens to sue for public fines payable to the Treasury raises ``difficult and fundamental questions'' under Article II's Take Care Clause and executive-branch authority, which remained open for future cases because the parties had not raised them.

\subsection{Dissenting Opinion (Justice Scalia)}

Justice Scalia attacked both halves of the majority's standing ruling. First, he argued that accepting subjective member affidavits about curtailed recreation despite the District Court's finding of zero actual environmental harm collapsed \textit{Lujan}'s injury-in-fact standard into a ``mere pleading requirement.'' Second, he contended that civil penalties paid exclusively to the U.S. Treasury cannot redress a private plaintiff's injury under Article III; any generalized deterrent effect is an undifferentiated public benefit, not personalized judicial relief. Echoing Article II concerns, he warned that the Court had placed the sovereign executive power to enforce public laws into unaccountable private hands.

\subsection{Subsequent Developments}

\begin{itemize}

\item \textbf{Enduring Authority on Voluntary Cessation (\textit{Already} and \textit{FBI v. Fikre}):} The Court repeatedly reaffirmed \textit{Laidlaw}'s ``formidable burden'' for voluntary cessation---first in \href{https://www.law.cornell.edu/supremecourt/text/11-982}{\textit{Already, LLC v. Nike, Inc.}, 568 U.S. 85 (2013)} (finding mootness only where Nike issued an unconditional, irrevocable covenant not to sue), and most prominently in \href{https://www.supremecourt.gov/opinions/23pdf/601us1r11_7lh8.pdf}{\textit{FBI v. Fikre}, 601 U.S. 234 (2024)} (\href{https://supreme.justia.com/cases/federal/us/601/22-1178/}{Justia}). In a unanimous opinion by Justice Gorsuch, \textit{Fikre} relied directly on \textit{Laidlaw} (528 U.S. at 189) to hold that the federal government's post-suit removal of a U.S. citizen from the No Fly List did not moot his due-process suit because a bare declaration promising not to relist him based on ``currently available information'' failed to prove the challenged conduct could not reasonably be expected to recur.

\item \textbf{Tension With \textit{TransUnion} on Concrete Injury:} By contrast, \textit{Laidlaw}'s expansive view of injury and citizen-suit redressability stands in growing tension with \href{https://www.law.cornell.edu/supremecourt/text/20-297}{\textit{TransUnion LLC v. Ramirez}, 594 U.S. 413 (2021)} (\href{https://supreme.justia.com/cases/federal/us/594/20-297/}{Justia}). Echoing Justice Scalia's \textit{Laidlaw} dissent, \textit{TransUnion} held that Congress cannot create Article III injury merely by enacting a statutory cause of action; plaintiffs must show concrete harm bearing a ``close historical or common-law analogue,'' and unmaterialized risk of future harm cannot support retrospective damages.

\end{itemize}

\section{Examples: Future Applications}
\begin{enumerate}

\item \textbf{Same-Side (Live Case; Not Moot): Post-Filing Compliance by a Chemical Plant.} A state environmental group whose members regularly kayak on Buffalo Bayou sues a chemical plant under CWA \S~505 while the plant is actively exceeding its benzene permit limits, seeking civil penalties payable to the Treasury. Six months into discovery, the plant installs new scrubbers, achieves full compliance, and moves to dismiss as moot. Under \textit{Laidlaw} and \textit{FBI v. Fikre}, the group had standing at filing (curtailed recreational use plus ongoing violations redressable by deterrent civil penalties), and the plant's post-filing compliance does not moot the case unless the plant proves it is ``absolutely clear'' that benzene violations cannot reasonably be expected to recur.

\item \textbf{Same-Side (Standing Exists): Birdwatchers Who Stay Away From a Marsh.} A paper mill upstream of a coastal marsh repeatedly exceeds the dioxin limits in its discharge permit. Members of a birding club submit affidavits that they used to visit the marsh monthly but have stopped because they fear the illegal discharges, and one says her lakeside cabin has become harder to rent. The mill shows that scientists have found no measurable harm to the marsh's wildlife. Under \textit{Laidlaw}, the relevant injury is to the plaintiffs, not the environment: a reasonable concern about ongoing unlawful pollution that leads members to give up specific, regular uses of an identified place is concrete injury in fact, unlike the vague some-day intentions in \textit{Lujan}.

\item \textbf{Opposite-Side (No Standing): Violations Ended Before Suit.} Suppose the same chemical plant cured all benzene violations and permanently surrendered its discharge permit \textit{before} the environmental group filed suit. Standing is measured when the complaint is filed, and at that moment nothing remained to deter. Civil penalties payable to the Treasury for wholly past violations cannot redress the members' injury, which is why \textit{Steel Co.\ v.\ Citizens for a Better Environment} found no standing. The voluntary-cessation doctrine never comes into play, because it applies only to a case that was live when filed and is later claimed to be moot.

\item \textbf{Opposite-Side (No Standing): An Undisclosed Credit-File Error.} Alternatively, suppose a consumer sues a credit bureau for statutory damages over an internal file error that was never disclosed to any third party. The consumer has a statutory cause of action, but under \textit{TransUnion} a bare statutory violation is not enough: the inaccurate file, never sent to anyone, causes no concrete harm with a close historical analogue such as defamation. Unlike the \textit{Laidlaw} members, who changed how they used a real place because of ongoing pollution, the consumer has suffered nothing Article III recognizes, so the suit must be dismissed.

\item \textbf{Fence-Sitter (Close Call): Closing the Plant but Keeping the Permit.} After a citizen suit is filed over ongoing permit violations, the defendant shuts down its facility, lays off its workers, and lists the property for sale, but keeps its discharge permit in force and says it has no present plans to reopen. That is essentially Laidlaw's own posture, which the Court remanded without deciding. The defendant argues that closure makes recurrence impossible; the plaintiffs answer that a retained permit leaves the defendant free to resume discharges at will, and \textit{Fikre} held that a bare assurance of no present intent does not carry the formidable burden. Whether the closure is truly permanent is a factual question that could go either way.

\end{enumerate}

\section{Critique}
\textbf{Reasonable Fear vs. Objective Harm:} Critics echo Justice Scalia's concern that separating ``injury to the plaintiff'' from ``injury to the environment'' risks manufacturing standing from subjective anxiety whenever a defendant violates a technical regulatory threshold, even when the violation causes zero physical or ecological harm.

\textbf{The Unresolved Article II Problem:} Justice Kennedy's concurrence flagged a structural tension that \textit{Laidlaw} left unanswered: when private plaintiffs sue for public Treasury fines without executive supervision, citizen suits may trench on the President's Article II prosecutorial discretion under the Take Care Clause.

\textbf{Standing--Mootness Asymmetry:} \textit{Laidlaw} decisively refutes the simplistic exam shorthand that ``mootness is just standing set in a time frame.'' Always remember that burden allocation shifts from plaintiff (standing) to defendant (mootness), and voluntary cessation saves a live suit that would have failed standing if filed after compliance.

\section{Key Quotations}
\begin{quote}
``The relevant showing for purposes of Article III standing \ldots{} is not injury to the environment but injury to the plaintiff.''
\end{quote}
\noindent --- Justice Ginsburg, for the Court, 528 U.S. at 181 (\href{https://www.law.cornell.edu/supct/html/98-822.ZD.html}{Cornell LII}; \href{https://supreme.justia.com/cases/federal/us/528/167/}{Justia}).

\begin{quote}
``To the extent that they encourage defendants to discontinue current violations and deter them from committing future ones, they afford redress to citizen plaintiffs who are injured or threatened with injury as a consequence of ongoing unlawful conduct.''
\end{quote}
\noindent --- Justice Ginsburg, for the Court, 528 U.S. at 185--186 (\href{https://www.law.cornell.edu/constitution-conan/article-3/section-2/clause-1/redressability}{Cornell LII ConAn}; \href{https://supreme.justia.com/cases/federal/us/528/167/}{Justia}).

\begin{quote}
``[A] defendant claiming that its voluntary compliance moots a case bears the formidable burden of showing that it is absolutely clear the allegedly wrongful behavior could not reasonably be expected to recur.''
\end{quote}
\noindent --- Justice Ginsburg, for the Court, 528 U.S. at 190 (\href{https://www.law.cornell.edu/supremecourt/text/11-982}{Cornell LII (quoting \textit{Laidlaw})}; \href{https://www.law.cornell.edu/supremecourt/text/528/167}{Cornell LII}).

\begin{quote}
``By uncritically accepting vague claims of injury, the Court has turned the Article III requirement of injury in fact into a `mere pleading requirement,' \ldots{} and by approving the novel theory that public penalties can redress anticipated private wrongs, it has come close to `mak[ing] the redressability requirement vanish.' \ldots{} The undesirable and unconstitutional consequence of today's decision is to place the immense power of suing to enforce the public laws in private hands.''
\end{quote}
\noindent --- Justice Scalia, dissenting, 528 U.S. at 215 (\href{https://www.law.cornell.edu/supct/html/98-822.ZD.html}{Cornell LII}).


\end{document}
